% !PRINTED 371
\weberpage{371}
From this it follows, with regard to
\[
\left(\frac{b-\sqrt D}{2}\right)^2 = a_1 a_2 c + b\,\frac{b-\sqrt D}{2}:
\]
\weq{11}{
\mu\nu = a_1 a_2 (u_1 v_1 - c\,u_2 v_2) + \frac{b-\sqrt D}{2}\,(a_1 u_1 v_2 + a_2 u_2 v_1 + b\,u_2 v_2).
}
If, then, one makes the bilinear substitution:
\weq{12}{
\begin{aligned}
t_1 &= u_1 v_1 - c\,u_2 v_2, \\
t_2 &= a_1 u_1 v_2 + a_2 u_2 v_1 + b\,u_2 v_2,
\end{aligned}
}
there results
\weq{13}{\lambda = \mu\nu,}
and consequently
\weq{14}{T = UV.}
In the requirement that $a_1$ and $a_2$ are to be relatively prime there lies at times a certain inconvenience, e.\,g. when it is a question of the composition of a form with itself.

Let us then drop the hypothesis that $\mathfrak{a}$ is relatively prime to $\mathfrak{b}$, and $a_1$ to $a_2$, but hold fast to the assumption that
\weq{15}{b^2 - D = 4\,a_1 a_2 c}
is divisible by $4 a_1 a_2$; then relation (11) still holds, and it follows that $\mathfrak{a}\mathfrak{b}$ is the greatest common divisor of
\[
a_1 a_2, \quad a_1\,\frac{b-\sqrt D}{2}, \quad a_2\,\frac{b-\sqrt D}{2}, \quad b\,\frac{b-\sqrt D}{2}
\]
(Bd.\,II, §~160, 2).

If now $a_1$, $a_2$, $b$ have no common divisor, then the equation $k_1 a_1 + k_2 a_2 + kb = 1$ can be satisfied by rational integers $k$, and accordingly $\mathfrak{a}\mathfrak{b}$ is likewise the greatest common divisor of
\weq{16}{a_1 a_2, \quad \frac{b-\sqrt D}{2}.}

2. \gesperrt{It is to be proved that $a_1 a_2$ is the smallest natural number divisible by $\mathfrak{a}\mathfrak{b}$.}

For if we set this smallest number $=a$, then $a_1 a_2$ is divisible by $a$, say
\[
a_1 a_2 = a\,a'.
\]
